\section{Perception}

The inward turn begins with perception, the simplest thing a self does with its world. A dock takes in the tones that stream to it, its fill, and sends its own back out, its will, and that exchange is the whole of the contact. To perceive is to let the incoming tones shape the inner ones, under a rule that conserves charge, so nothing is taken in without something given back. Perception here is not a window onto an outside. It is the meeting of two flows at a boundary, the self's tones and the world's, and what the self feels of the world is the trace the world leaves as the two settle. The richer the boundary, the more of the world it can hold, which is why a mind perceives more than a stone, not by looking harder but by having more inside for the world to write on.
